\section{Mistral 7B: data-first y diseño de modelos compactos}

La historia que se cuenta en clase sobre Mistral 7B retoma dos juicios: primero, la calidad de los datos es la palanca en la que más conviene invertir para un equipo pequeño; segundo, un modelo pequeño solo convierte el hecho de «tener pocos parámetros» en una ventaja real de despliegue si optimiza a la vez el entrenamiento y la arquitectura de inferencia. El artículo oficial de Mistral 7B presenta dos mecanismos centrales: Grouped-Query Attention (GQA) y Sliding Window Attention (SWA).

\subsection{Qué problemas resuelven GQA y SWA}

\begin{knowledgebox}{Términos clave: mecanismos de atención de los modelos compactos}
\begin{tabularx}{\textwidth}{L{0.18\textwidth}>{\raggedright\arraybackslash}X>{\raggedright\arraybackslash}X}
\toprule
Término & Problema que resuelve & Mecanismo central y relación con el despliegue \\
\midrule
GQA & KV cache y ancho de banda de decodificación demasiado grandes & Varias query heads comparten un número menor de key/value heads, lo que reduce la KV cache y el volumen de lectura, y aumenta el inference throughput. \\
SWA & Complejidad cuadrática de la full attention en contextos largos & Cada capa atiende directamente solo una ventana local; la información se propaga gradualmente a través de varias capas, lo que amplía el campo receptivo efectivo a menor costo. \\
KV cache & La decodificación autorregresiva recalcula las key/value del historial & Guarda los estados de atención del historial para evitar recalcularlos; su capacidad suele ser la principal restricción para el batch size y la concurrencia. \\
\bottomrule
\end{tabularx}
\end{knowledgebox}

El artículo de Mistral 7B también publica los pesos y una implementación de referencia bajo Apache 2.0. El significado sistémico aquí no es que «la licencia crea un ecosistema por sí sola», sino que permite a los equipos desplegar, modificar y hacer fine-tune del modelo de forma local o en la nube, de modo que el modelo base se trata como un activo componible. Los pesos abiertos amplían la experimentation surface, pero también trasladan en mayor medida al integrador la responsabilidad de la seguridad, la evaluación y la operación.

\begin{importantbox}{La ventaja de un modelo pequeño debe reflejarse en cuatro dimensiones de recursos}
\begin{enumerate}
  \item \textbf{Memoria de GPU}: si los parámetros, la KV cache y el runtime workspace caben en el hardware de destino;
  \item \textbf{Ancho de banda}: cuántos pesos y cuánta cache hay que leer de la HBM por cada token;
  \item \textbf{Latencia}: si el usuario puede aceptar la latencia del primer token y la latencia entre tokens;
  \item \textbf{Rendimiento}: cuántas solicitudes concurrentes se pueden atender con un nivel de servicio dado.
\end{enumerate}
Si la ventaja en un Benchmark no mejora este balance de recursos, no necesariamente se traduce en una ventaja de despliegue.
\end{importantbox}

\begin{warningbox}{Un artículo de arquitectura no sustituye a un workload benchmark}
GQA y SWA ofrecen ventajas a nivel de mecanismo, pero el beneficio real depende de la longitud de secuencia, el batching, la implementación de los kernels, la cuantización, el hardware y los objetivos de servicio. Al elegir un modelo, hay que medir con el tráfico y los datos de destino, en lugar de limitarse a copiar la puntuación promedio del artículo.
\end{warningbox}

\subsection{Resumen de la sección}

El valor didáctico de Mistral 7B está en que conecta una forma de organización data-first con una arquitectura inference-aware. Un modelo compacto no es una versión reducida de un modelo grande, sino una redistribución del presupuesto de entrenamiento y de arquitectura en torno a las restricciones de recursos de un despliegue viable.
